%% ============================================================
%% CHAPTER 1 — INTRODUCTION
%% ============================================================

\chapter{Introduction}
\label{ch:introduction}

\section{Motivation and Strategic Context}
\label{sec:intro_motivation}

Natural gas transmission pipelines are load-bearing arteries of modern
states: the \ac{TEN-E} network alone delivers \SI{>400}{bcm}
per year across 27 EU member states, while Gulf Coast corridors in the
United States carry a comparable volume domestically~\cite{PHMSA2024}.
Their disruption — whether by material failure, physical sabotage, or
cyber-physical manipulation — constitutes a dual economic and security
crisis.  In the European context, this is not a theoretical concern.
The 2022 Nord Stream I/II ruptures — the largest deliberate destruction
of subsea gas infrastructure in history — destroyed approximately
\SI{1200}{\kilo\meter} of DSAW-welded pipe operating at \SI{22}{\mega\pascal}
and removed \SI{55}{bcm\,yr^{-1}} of transmission
capacity from European markets overnight.  In the United States, the
2021 Colonial Pipeline ransomware incident shut \SI{8850}{\kilo\meter}
of refined-product pipeline for six days, triggering regional fuel
shortages across the southeastern seaboard and a declaration of emergency
in 17 states.  Both events illustrate that pipeline infrastructure sits
at the precise intersection of structural engineering, network resilience,
and adversarial intent — domains that have, until now, been modelled
entirely in isolation.

The strategic importance of pipeline security is underscored by the
legislative response.  The EU Critical Entities Resilience Directive
(CER, 2022/2557) explicitly identifies gas transmission networks as
critical entities and mandates risk assessments that account for both
physical sabotage and cyber-physical threats by 2026.  Mediterranean
corridors — the Trans-Adriatic Pipeline (\ac{TAP}), the Trans-Anatolian
Pipeline (TANAP), and the SNAM~Rete~Gas onshore network — are classified
as Projects of Common Interest under the revised \ac{TEN-E} Regulation
(EU~2022/869) and carry particular strategic weight as diversification
assets in the post-Nord~Stream security environment.  Any analytical
framework that aspires to inform resource allocation decisions for these
assets must therefore bridge structural integrity and adversarial
strategy simultaneously.

Contemporary risk frameworks fail this test.  Regulatory standards such
as \ac{ASME}~B31.8~\cite{ASME_B31_8:2025} and \ac{BS7910}:2019~\cite{BS7910:2019}
provide rigorous tools for assessing structural integrity but make no
provision for strategic adversaries — they quantify the probability that
a flawed weld fails under pressure loading, not the probability that a
rational actor will target that weld.  Security game models have
demonstrated impressive real-world deployments — from \ac{LAX} patrol
scheduling~\cite{An2013} to ferry route randomisation~\cite{Tambe2011} —
but universally assign attacker payoffs from abstract scalar constants
supplied by subject-matter experts, entirely uncoupled from physical
reality.  The failure probability $\Pf = 0.88$ of a fillet-welded
legacy segment at maximum allowable operating pressure is categorically
different from $\Pf = 0.53$ for a seamless X65 section; yet no extant
Bayesian Stackelberg formulation encodes this distinction.  A technically
sophisticated attacker who reads BS~7910 enjoys a systematic information
advantage over a defender whose game model ignores the standard entirely.

This thesis proposes and validates a remedy: a physics-informed payoff
function derived from first-principles fracture mechanics, calibrated against
\num{1996} \ac{PHMSA} gas transmission incidents (2010--2024), and embedded
directly into a Bayesian \ac{SSE}~\cite{Paruchuri2008} over a realistic
pipeline network topology.  A second, equally novel contribution addresses
the \emph{adversarial NDE threat}: the possibility that a sophisticated
actor could manipulate the very inspection system designed to detect
structural deterioration, forcing a physically vulnerable weld to register
as sound — an intelligence preparation of the operational environment that
creates the conditions for a subsequent physical attack.

\section{Problem Statement}
\label{sec:intro_problem}

Let $G = (\mathcal{V}, \mathcal{E})$ be a pipeline network where each
edge $e \in \mathcal{E}$ is characterised by physical attributes
$(D_e, t_e, \sigma_{y,e}, \text{seam}_e)$ — outer diameter, wall
thickness, yield strength, and seam type.  The \emph{physics-informed
attacker utility} is:
\begin{equation}
  \Ua(e) = w_{\Pf}\,\Pf(e) + w_v\,v(e) + w_b\,b(e),
  \label{eq:attacker_utility}
\end{equation}
where $\Pf(e)$ is the segment failure probability derived from BS~7910
Level~2 fracture assessment and \acs{IIW} fatigue analysis,
$v(e)$ is a normalised strategic value (diameter $\times$ length), and
$b(e)$ is the normalised betweenness centrality.  The weights
$w_{\Pf}, w_v, w_b$ are type-specific for each of the three adversary
profiles defined in \cref{ch:physics}.

The defender allocates a budget $B$ by choosing coverage probabilities
$\bmc \in [0,1]^{|\mathcal{E}|}$ to minimise network-level expected loss.
This yields the following three-part research question:

\begin{enumerate}
  \item Does replacing abstract payoffs with $\Pf$ values derived from
    \acs{BS7910} fracture mechanics produce a qualitatively different
    \ac{SSE} than uniform vulnerability assumptions?
  \item How does the physics-informed \ac{SSE} compare to uniform
    resource allocation in terms of total network expected loss?
  \item Can a white-box adversary suppress the $\Pf$ signal by
    manipulating \ac{NDE} inspection data at the segment level?
\end{enumerate}

\section{Thesis Contributions}
\label{sec:intro_contributions}

This thesis makes the following original contributions:

\begin{enumerate}
  \item \textbf{Physics-Informed Payoff Engine.}  A \ac{MC} failure
    probability model grounded in \ac{BS7910}:2019 Level~2 \ac{FAD}
    and \acs{IIW} fatigue classification, calibrated against 216
    material/weld failure events extracted from the \ac{PHMSA} incident
    database.  The engine reproduces the empirical failure-rate hierarchy
    across six seam types (seamless to fillet-welded legacy pipe) with
    $\Pf \in [\num{0.29}, \num{0.93}]$ across the test network
    (\cref{ch:physics}).

  \item \textbf{Bayesian Stackelberg Security Game.}  A three-attacker-type
    Bayesian \ac{SSE} with LP-enumeration solver, where attacker utilities
    are supplied directly by the physics engine.  At a \SI{30}{\percent}
    defender budget, the physics-informed allocation achieves
    $+\SI{52.3}{\percent}$ coverage effectiveness and a \SI{17.0}{\percent}
    reduction in total expected network loss versus a uniform baseline
    (\cref{ch:physics}).

  \item \textbf{Adversarial NDE Threat Model.}  A from-scratch NumPy
    \ac{MLP} weld defect classifier (32-feature, 99.7\% clean accuracy),
    attacked by \ac{FGSM}~\cite{Goodfellow2015},
    \ac{BIM}~\cite{Kurakin2017}, and \ac{PGD}~\cite{Madry2018} with
    physics-informed $\eps$ scaling.  \ac{BIM} achieves the highest
    \ac{ASR} of \SI{20.1}{\percent} at $\eps = \num{0.30}$, revealing
    the exploitable gap between iterative and random-start attacks on this
    feature geometry (\cref{ch:adversarial}).

  \item \textbf{Integrated Dashboard.}  A Plotly Dash application unifying
    all four model layers — network \ac{FAD}, Stackelberg coverage, \ac{NDE}
    adversarial impact — with click-to-inspect segment intelligence and
    scenario comparison (\cref{ch:results}).  The implementation comprises
    310 unit tests with full public-interface coverage.
\end{enumerate}

\section{Thesis Scope and Limitations}
\label{sec:intro_scope}

The network topology used throughout is \emph{synthetic}: a 20-node,
22-segment Gulf Coast transmission corridor generated to match \ac{PHMSA}
fleet statistics rather than any specific operator's infrastructure.  This
is a deliberate choice to avoid intellectual property constraints while
retaining physical realism.  The \ac{NDE} classifier is trained on
synthetic signal features rather than real radiograph images, as no
public dataset of labelled weld radiographs at the required scale is
available.  The adversarial attacks are therefore evaluated in a
feature-vector threat model (\ac{FGSM}/\ac{BIM}/\ac{PGD} on z-score
normalised features) rather than against raw image data.

Extensions to real network topologies (TAP, SNAM~Rete~Gas) and image-based
NDE pipelines using convolutional architectures are discussed in
\cref{ch:conclusion}.

\section{Regulatory Framework}
\label{sec:intro_regulatory}

The physics-informed failure assessment operates within a hierarchy of
international standards, each playing a distinct role in the analytical
chain.

\ac{BS7910}:2019~\cite{BS7910:2019} is the primary assessment standard,
providing the Level~2 Option~1 \ac{FAD} curve, the stress intensity
factor solution for surface flaws in cylindrical shells (Annex~M), and
the partial safety factor calibration used for probabilistic assessments
(Annex~K).  It is the preferred standard for European pipeline operators
and is the basis for all fracture assessments in \cref{ch:physics}.

\acs{ASME}~B31.8:2025~\cite{ASME_B31_8:2025} governs pressure design
and integrity management for gas transmission piping systems in the
United States, from which the \acs{PHMSA} incident database is drawn.
Its \SI{72}{\percent}~\acs{SMYS} design factor and high-consequence area
(HCA) classification are used in setting the operating pressure boundary
conditions for the Monte Carlo engine.

\acs{API}~1104:2021~\cite{API1104:2021} provides the acceptance criteria
for weld flaws detected during construction inspection and in-service
in-line inspection.  These criteria define the upper bound on flaw
dimensions used to truncate the Monte Carlo sampling distributions:
any flaw exceeding the API~1104 acceptance limit would trigger
immediate remediation and is therefore excluded from the long-term
failure probability assessment.

\acs{API}~579-1/\acs{ASME}~FFS-1:2021~\cite{API579:2021} serves as a
cross-reference for the reference stress solution and provides the
alternative FFS Level~2 acceptance criteria against which the
BS~7910 results are compared in \cref{subsec:physics_kr_lr}.  The
two standards agree within \SI{3}{\percent} on $L_r$ for the pipeline
geometry range used in this thesis.

The \acs{IIW} Fatigue Design Recommendations~\cite{IIW_Fatigue:2016}
provide the S-N curve classification system (FAT class assignments)
and the bilinear damage model used in the \texttt{fatigue\_engine.py}
module.  DNVGL-RP-F101:2017~\cite{DNVGL_RPF101:2017} is cited for
corrosion-related wall thickness input distributions.

\noindent All numerical values in \cref{ch:physics} that reference these
standards are drawn from the edition and clause specified above; no
earlier editions are used.

\section{Thesis Structure}
\label{sec:intro_structure}

The remainder of the thesis is organised as follows.
\Cref{ch:literature} surveys the three foundational bodies of literature:
pipeline structural integrity, Stackelberg security games, and adversarial
robustness of \ac{NDE} systems.
\Cref{ch:physics} develops the physics-informed vulnerability engine
(Zone~C): the \ac{MC} $\Pf$ model, the network graph, and the Bayesian
\ac{SSE}.
\Cref{ch:adversarial} presents the adversarial \ac{NDE} threat model
(Zone~A): the \ac{MLP} classifier, its training, and the three attack
implementations.
\Cref{ch:results} integrates all layers, presents the 17.0\% risk
reduction finding, and analyses the dashboard scenario comparison.
\Cref{ch:conclusion} synthesises the contributions, discusses applicability
to Mediterranean pipeline corridors, and proposes six directions for
future research.
